\documentclass[10pt,a4paper]{amsart}
\usepackage[margin=19mm]{geometry}
\usepackage{amsmath,amssymb,mathtools}
\usepackage[scr=rsfs,cal=euler]{mathalfa}
\usepackage{xfrac}
\usepackage[unicode,hidelinks,bookmarks=false]{hyperref}
\newcommand{\ie}{i.e.\ }
\newcommand{\cf}{cf.\ }
\newcommand{\resp}{respectively\ }
\newcommand{\Subsection}{\subsection}
\providecommand{\nobreakdash}{\nobreak\hbox{-}}
\setlength{\emergencystretch}{2em}
\multlinegap0pt
\usepackage{bm}
\usepackage{bbm}
\usepackage{dsfont}
\usepackage{xspace}
\usepackage{cancel}
\usepackage{mathtools}
\usepackage{amsthm}
\newtheorem{theorem}{Theorem}
\title[Local certification of geometric graph classes]{Local certification of geometric graph classes}
\author{Oscar Defrain, Louis Esperet, Aur\'elie Lagoutte, Pat Morin, Jean-Florent Raymond}
\date{Source: arXiv:2311.16953v5 (CC BY 4.0)}
\begin{document}
\maketitle

\noindent\emph{Source and licence.} Local certification of geometric graph classes. Version arXiv:2311.16953v5 (CC BY 4.0), distributed under the Creative Commons Attribution 4.0 International licence, as recorded in the arXiv licence metadata for this exact version and shown on the abstract page https://arxiv.org/abs/2311.16953v5. Full rights notice: licenses/arxiv-2311.16953-CC-BY-4.0.txt.\par\smallskip

\noindent\textit{Editorial scope.} Counted unit: the Theorem whose source label is \texttt{None} in the article (source0.tex lines 590--594, source label \texttt{None}), delivered together with the article's own complete original proof of it (source0.tex lines 596--642, one \texttt{proof} environment of 47 lines). No mathematics is written, repaired or re-derived: statement and proof are the article's own text line for line. Required context retained uncounted: none. Every definition, display and label of those lines is retained unchanged. Omitted from this package and not redistributed: 1--589, 595--595, 643--2516. Nothing inside a retained excerpt is omitted or altered: every retained body is delivered exactly as the article's lines are written, so no omission receipt of any kind accompanies this package. Citations: the retained excerpts carry no citation command at all, so no bibliography entry of the article is transcribed and no reference is redistributed. Reference adaptation: none was needed and none was made; every reference inside the delivered unit resolves inside it, so no unresolved reference or citation is produced. Printed slips kept literal: no printed word, symbol, hypothesis or number of the counted statement or of its proof was altered. Layout and class: the article's own class is \texttt{amsart} with packages including geometry, inputenc, fontenc, amsthm, amssymb, amsmath, thmtools, thm-restate, float, needspace, graphicx, xcolor, subcaption, hyperref; the wrapper is the lane's pinned \texttt{amsart} editorial wrapper at 10pt on A4, which restates the theorem-like environments the retained text needs and adds the article's own macro definitions that the retained text needs (none), with extra wrapper packages (bm, bbm, dsfont, xspace, cancel, mathtools). The delivered numbering is the unit's own. Assets: no retained span contains a figure environment, a graphics-inclusion command, a PDF-page inclusion, a table or a TikZ picture; no image file is copied into this package, the package declares no asset dependency and no image file is redistributed with it. Licence: version arXiv:2311.16953v5 of this article is distributed under the Creative Commons Attribution 4.0 International licence, as recorded in the arXiv licence metadata for this exact version and shown on the abstract page https://arxiv.org/abs/2311.16953v5. A full rights notice is delivered at licenses/arxiv-2311.16953-CC-BY-4.0.txt. Characters: every retained excerpt is pure ASCII, so the delivered engine, XeLaTeX with its default Latin Modern text font, reports no missing character.
\begin{theorem}
  Any class $\mathcal{C}$ of connected graphs has local complexity at most
$\log(|\mathcal{C}_n|)+O(\log n)$. In particular if  $\mathcal{C}$ is
a tiny class, then the local complexity is $O(n)$.
\end{theorem}

\begin{proof}
Let $G\in \mathcal{C}_n$. The certificate given by the prover to
each vertex $v$ of $G$ contain the following:

\begin{itemize}
\item the number of vertices $n$;
\item a $\log(|\mathcal{C}_n|)$-bit word $w$ representing a graph $G'$ isomorphic to $G$;
\item the name of the vertex $\pi(v)$ corresponding to $v$ in $G'$.
\end{itemize}
In addition to this, the vertices of $G$ store a locally certified spanning
tree $T$, rooted in some vertex $r\in G$, which can be done with
$O(\log n)$ additional bits per vertex (this also encodes the
parent-child relation in the tree $T$, so that each vertex knows that
it is the root $r$ or knows the identifier of its parent in the
tree). We also give to each vertex $v$ the number of vertices in its
rooted subtree $T_v$.
In total, the certificates
above take $\log(|\mathcal{C}_n|)+O(\log n)$ bits, as desired.

\medskip

We now describe the verifier part. Each
vertex checks that it has been given the same value of $n$ and the same
word $w$ describing some graph $G'\in \mathcal{C}$ as its neighbors. The spanning tree $T$ is
then used to compute the number of vertices of $G$ and the root $r$
checks that this number coincides with $n$ and the number of vertices
of $G'$ (this is standard: each
vertex $v$ checks that the number of vertices in its rooted subtree
$T_v$, which was given as a certificate, is equal to the sum of the
number of vertices in the rooted subtrees of its children in $T$, plus
1).
Then each vertex $v$
verifies that $\pi$ is a local isomorphism from $G$ to $G'$, that is,
$\pi$ maps bijectively the neighborhood of $v$ in $G$ to the neighborhood of
$\pi(v)$ in $G'$.

\medskip

We now analyze the scheme. If $G\in \mathcal{C}$, then clearly all the
vertices accept. Assume now that all the vertices accept. Then $\pi$
is a local isomorphism from $G$ to some graph $G'\in \mathcal{C}$,
with the same number of vertices as $G$. As $G'$ is connected, $\pi$
is surjective, but as $G$
and $G'$ have the same number of vertices, $\pi$ must also be
injective. Thus $\pi$ is a bijection and $G$ and $G'$ are isomorphic,
which implies that $G\in \mathcal{C}$.
\end{proof}

\end{document}
